\subsection{Icon changes}
The previous development group working on this product \citep{misc:crocold} performed usability tests and through these, found that some of the icons was not intuitive to the users and thus needed changing.
When changing the icons, the focus was to make them easier to associate the functionality to the icons.
The idea behind the new choice of freehand drawing icon is what people associate with this functionality and how it is represented in similar programs.
In such similar programs the standard representation of a freehand drawing is a pencil.
Due to this standard the icon in \figref{figure:old-freehand} was replaced by \figref{figure:new-freehand}.

\begin{figure}[h]
	\centering
	%---- linebreak	
	\begin{subfigure}[b]{0.45\textwidth}
		\centering
		\includegraphics[scale = 1]{media/freehandNew}
		\caption{New freehand drawing icon}
		\label{figure:new-freehand}
	\end{subfigure}
	\qquad
	\begin{subfigure}[b]{0.45\textwidth}
		\centering
		\includegraphics[scale = 0.1]{media/freehandOld}
		\caption{Old freehand drawing icon}
		\label{figure:old-freehand}
	\end{subfigure}
	\caption{Freehand drawing icons}
	\label{figure:freehand}
\end{figure}

The previous icon for the selection tool could give a misleading impression of the functionality.
The icon, seen in \figref{figure:old-select}, looks much like a tool that marks an area, but the tool only selected one object.
To make the icon represent the functionality of the feature a new icon was deemed necessary.
The new icon was made to look like a hand, as this gives a clear representation of only being able to select one object, see \figref{figure:new-select}.

\begin{figure}[h]
	\begin{subfigure}[b]{0.45\textwidth}
		\centering
		\includegraphics[scale = 1]{media/selectNew}
		\caption{New selection icon}
		\label{figure:new-select}
	\end{subfigure}
	\qquad
	\begin{subfigure}[b]{0.45\textwidth}
		\centering
		\includegraphics[scale = 0.1]{media/selectOld}
		\caption{Old selection icon}
		\label{figure:old-select}
	\end{subfigure}		
	\caption{Selection icons}
	\label{figure:select}
\end{figure}